
\section{\BookRecordIndex{T0099}等周不等式}\label{sec:4-4}
\zhindex{dengzhoubudengshi}{等周不等式}

最基本的\BookRecordIndex{T0099}等周不等式讨论欧氏平面中的简单闭曲线，通常表述为：在所有给定长度的此类
曲线中，圆所围成的面积最大。也可以这样说：若 $\gamma$ 是 $\mathbb R^2$ 中的简单
闭曲线，以 $L(\gamma)$ 表示它关于通常度量的长度，以 $A(\gamma)$ 表示它围成的面积，
则
\[
  A(\gamma)\le\frac1{4\pi}L(\gamma)^2,
\]
且等号成立当且仅当 $\gamma$ 是圆。本节将说明，这一形式的\BookRecordIndex{T0099}等周不等式可以自然推广到
\BookRecordIndex{T0358}辛流形；其中 $A(\gamma)$ 解释为以 $\gamma$ 为边界的圆盘之辛面积。一般而言（即
$\omega$ 在 $\pi_2(M)$ 上不消失时），$\gamma$ 是许多辛面积不同的圆盘之边界，所以
必须指明所考察的圆盘。只有当 $\gamma$ 充分短时，才能做到这一点。

下面的定理~\ref{thm:4-4-1} 在 $\omega$-\BookRecordIndex{T0348}相容情形以最佳常数建立\BookRecordIndex{T0099}等周不等式。利用它，
可以在引理~\ref{lem:4-7-3} 的指数估计中，对 $\omega$-\BookRecordIndex{T0348}相容\BookRecordIndex{T0088}殆复结构得到最佳常数。
$\omega$-\BookRecordIndex{T0372}驯服情形的一般\BookRecordIndex{T0099}等周不等式（不控制常数）则可用较容易的方法证明，见
注~\ref{rem:4-4-3} 的概述；证明奇点可去性只需要这个一般形式。

设 $(M,\omega)$ 是紧\BookRecordIndex{T0358}辛流形，$J\in\mathcal J_\tau(M,\omega)$ 是被 $\omega$ \BookRecordIndex{T0372}驯服的
\BookRecordIndex{T0088}殆复结构。以下始终给 $M$ 配备由 $\omega$ 和 $J$ 通过式~\eqref{eq:2-1-1} 确定的黎曼度量
$g_J=\langle\cdot,\cdot\rangle_J$。考虑光滑闭路
$\gamma:\mathbb R/2\pi\mathbb Z\to M$，并假设其长度
\[
  \ell(\gamma):=\int_0^{2\pi}|\dot\gamma(\theta)|_J\,d\theta
\]
小于 $M$ 的\BookRecordIndex{T0093}单射半径。于是它的像包含在某个测地球 $U\subset M$ 中，且该球的半径至多
为\BookRecordIndex{T0093}单射半径的一半。由于这个球可缩，$\gamma$ 有一个光滑的局部延拓
$u_\gamma:\mathbb D\to U$，使得对每个 $\theta\in\mathbb R$，
\[
  u_\gamma(e^{i\theta})=\gamma(\theta).
\]
把 $\gamma$ 的\emph{\BookRecordIndex{T0202}局部辛作用量}定义为 $u_\gamma$ 的辛面积之负值\footnote{这里的
负号看起来略显别扭；不过，它与辛作用量的标准符号约定一致。}，记作
\zhindex{jubuxinzuoyongliang}{局部辛作用量}
\[
  a(\gamma):=-\int_{\mathbb D}u_\gamma^*\omega.
\]
任意两个这样的延拓（其像均落在半径至多为\BookRecordIndex{T0093}单射半径一半的测地球中）都相对于边界同伦，
因而具有相同的辛面积。所以，只要 $\gamma$ 充分短，\BookRecordIndex{T0202}局部辛作用量 $a(\gamma)$ 就有良好
定义。还要注意，$\omega$ 在 $U$ 上的限制是正合的，因此也可以用下式定义\BookRecordIndex{T0202}局部辛作用量：
\begin{equation}\label{eq:4-4-1}
  a(\gamma):=\int_\gamma\lambda,
  \qquad
  \lambda\in\Omega^1(U),
  \qquad d\lambda=-\omega.
\end{equation}
这个定义与前一定义一致，由 Stokes 定理即知。

现在设 $L\subset M$ 是紧\BookRecordIndex{T0219}拉格朗日子流形。\BookRecordIndex{T0202}局部辛作用量的概念可如下推广到边界点位于
$L$ 中的道路。设 $\gamma:[0,\pi]\to M$ 是光滑道路，并且
$\gamma(0),\gamma(\pi)\in L$。若 $\gamma$ 充分短，则它的像落在一个可缩开集
$U\subset M$ 中，而且 $L\cap U$ 也可缩。因此存在半圆盘
$u_\gamma:\mathbb D\cap\mathbb H\to M$，使得
\[
  u_\gamma(\mathbb D\cap\mathbb R)\subset L,
  \qquad
  u_\gamma(e^{i\theta})=\gamma(\theta),
  \quad 0\le\theta\le\pi.
\]
$\gamma$ 的\BookRecordIndex{T0202}局部辛作用量定义为
\[
  a(\gamma):=-\int_{\mathbb D\cap\mathbb H}u_\gamma^*\omega.
\]
或者，也可用式~\eqref{eq:4-4-1} 定义 $a(\gamma)$，但须要求 $\lambda$ 在 $L$ 上消失。
由 Stokes 定理，这两个定义仍然一致。因此，\BookRecordIndex{T0202}局部辛作用量与定义中所用的映射
$u_\gamma$ 和 $1$-形式 $\lambda$ 无关。

\begin{theorem}[\BookRecordIndex{T0099}等周不等式]\label{thm:4-4-1}
设 $(M,\omega)$ 是紧\BookRecordIndex{T0358}辛流形，$L\subset M$ 是紧\BookRecordIndex{T0219}拉格朗日子流形，并设 $J$ 是 $M$ 上的
\BookRecordIndex{T0088}殆复结构。则有：
\begin{enumerate}[label=(\roman*),leftmargin=3em]
\item 若 $J$ 与 $\omega$ \BookRecordIndex{T0348}相容，则对每个常数 $c>1/(4\pi)$，都存在常数
$\delta>0$，使每条光滑闭路 $\gamma:S^1\to M$ 都满足
\begin{equation}\label{eq:4-4-2}
  \ell(\gamma)<\delta
  \quad\Longrightarrow\quad
  |a(\gamma)|\le c\ell(\gamma)^2.
\end{equation}
\item 若 $J$ 与 $\omega$ \BookRecordIndex{T0348}相容，则对每个常数 $c>1/(2\pi)$，都存在常数
$\delta>0$，使式~\eqref{eq:4-4-2} 对每条端点位于 $L$ 中的光滑道路
$\gamma:[0,\pi]\to M$ 成立。
\item 若 $J$ 被 $\omega$ \BookRecordIndex{T0372}驯服，则存在正常数 $c,\delta$，使式~\eqref{eq:4-4-2}
对每条光滑闭路 $\gamma:S^1\to M$，以及每条端点位于 $L$ 中的光滑道路
$\gamma:[0,\pi]\to M$ 成立。
\end{enumerate}
\end{theorem}

\begin{remark}\label{rem:4-4-2}
固定 $M$ 上一个与 $\omega$ \BookRecordIndex{T0348}相容的\BookRecordIndex{T0088}殆复结构 $J_0$。以 $\ell_0(\gamma)$ 表示闭路
$\gamma$ 关于度量 $g_{J_0}$ 的长度，并令 $\delta_0$ 为 $g_{J_0}$ 的\BookRecordIndex{T0093}单射半径的一半。
那么对每条满足 $\ell_0(\gamma)<\delta_0$ 的闭路，\BookRecordIndex{T0202}局部辛作用量 $a(\gamma)$ 都有良好
定义。给定 $J\in\mathcal J_\tau(M,\omega)$ 和闭路 $\gamma$，以 $\ell(\gamma)$ 表示
它关于 $g_J$ 的长度。令 $c_J\ge1$ 为使下式对每条闭路 $\gamma$ 成立的最小常数：
\[
  \frac1{c_J}\ell_0(\gamma)
  \le\ell(\gamma)\le c_J\ell_0(\gamma).
\]
定理~\ref{thm:4-4-1} 表明，常数 $\delta$ 只依赖于 $c_J$。所以，
定理~\ref{thm:4-4-1} 中的常数 $\delta=\delta(J)>0$ 可以选得关于
$\mathcal J_\tau(M,\omega)$ 上的 $C^0$ 拓扑连续。
\end{remark}

\begin{remark}\label{rem:4-4-3}
容易看出，\BookRecordIndex{T0099}等周不等式以某个常数 $c$ 成立。把 $\gamma$ 的局部延拓
$u_\gamma:\mathbb D\to M$ 定义为
\begin{equation}\label{eq:4-4-3}
  u_\gamma(re^{i\theta}):=\exp_{\gamma(0)}(r\xi(\theta)),
\end{equation}
其中 $\xi(\theta)\in T_{\gamma(0)}M$ 由条件
\[
  \exp_{\gamma(0)}(\xi(\theta))=\gamma(\theta)
\]
确定。为证明式~\eqref{eq:4-4-2}，注意
\[
  |\partial_ru_\gamma|
  =|\xi(\theta)|
  =d\bigl(\gamma(0),\gamma(\theta)\bigr)
  \le\ell(\gamma),
\]
而且存在只依赖于度量的常数 $c_1,c_2$，使得对所有 $r,\theta$，
\[
  |\partial_\theta u_\gamma|
  \le c_1|\dot\xi(\theta)|
  \le c_2|\dot\gamma(\theta)|.
\]
因此，应取\BookRecordIndex{T0202}局部辛作用量的绝对值来估计\sourcefootnote{原书此处写作
$|a(u_\gamma)|$；但 $a$ 在本节中只定义在边界闭路 $\gamma$ 上，而 $u_\gamma$ 是填充
圆盘，故据作用量的参数类型改作 $|a(\gamma)|$。}，并且只要 $\gamma$ 的长度小于
\BookRecordIndex{T0093}单射半径的一半，就有
\[
  |a(\gamma)|
  =\left|\int_0^{2\pi}\int_0^1
    \omega(\partial_ru_\gamma,\partial_\theta u_\gamma)
    \,dr\,d\theta\right|
  \le c_3\ell(\gamma)^2.
\]
\end{remark}

定理~\ref{thm:4-4-1} 的证明思路是：先直接计算，证明它在辛向量空间中成立；再利用
\BookRecordIndex{T0081}达布坐标卡\BookRecordIndex{T0326}推出它在一般\BookRecordIndex{T0358}辛流形上成立。这个办法之所以有效，是因为坐标卡不扭曲
作用量，而对短曲线的长度只产生很小的扭曲。

下面给出细节。设 $(V,\omega)$ 是辛向量空间，
$J\in\operatorname{Aut}(V)$ 是与 $\omega$ \BookRecordIndex{T0348}相容的复结构。于是 $V$ 配备了如
式~\eqref{eq:2-1-1}
所定义的内积
\[
  \langle v,w\rangle:=\omega(v,Jw).
\]
以下把单位圆 $S^1\subset\mathbb C$ 与 $\mathbb R/2\pi\mathbb Z$ 等同。闭路
$\gamma:S^1\to V$ 的作用量、长度与能量分别为
\[
  A(\gamma):=\frac12\int_0^{2\pi}
  \omega(\dot\gamma(\theta),\gamma(\theta))\,d\theta,
\]
\[
  L(\gamma):=\int_0^{2\pi}|\dot\gamma(\theta)|\,d\theta,
  \qquad
  E(\gamma):=\frac12\int_0^{2\pi}|\dot\gamma(\theta)|^2\,d\theta.
\]
把 $2\pi$ 换成 $\pi$，同样的公式定义光滑道路 $\gamma:[0,\pi]\to V$ 的作用量、
长度与能量。下一个引理以最佳常数建立辛向量空间中的\BookRecordIndex{T0099}等周不等式。

\begin{lemma}\label{lem:4-4-4}
设 $(V,\omega)$ 是辛向量空间，$\Lambda\subset V$ 是拉格朗日子空间，且
$J\in\operatorname{Aut}(V)$ 是 $V$ 上与 $\omega$ \BookRecordIndex{T0348}相容的复结构。则对每条光滑闭路
$\gamma:S^1\to V$，
\begin{equation}\label{eq:4-4-4}
  |A(\gamma)|\le\frac1{4\pi}L(\gamma)^2,
\end{equation}
而对每条满足 $\gamma(0),\gamma(\pi)\in\Lambda$ 的光滑道路
$\gamma:[0,\pi]\to V$，
\begin{equation}\label{eq:4-4-5}
  |A(\gamma)|\le\frac1{2\pi}L(\gamma)^2.
\end{equation}
\end{lemma}

\begin{proof}
先证明对每条闭路 $\gamma:S^1\to V$ 都有估计
\begin{equation}\label{eq:4-4-6}
  |A(\gamma)|\le E(\gamma).
\end{equation}
证明依赖恒等式
\begin{equation}\label{eq:4-4-7}
  \int_0^{2\pi}\omega(e^{k\theta J}v,e^{\ell\theta J}Jw)\,d\theta
  =
  \begin{cases}
    2\pi\langle v,w\rangle,&k=\ell,\\
    0,&k\ne\ell.
  \end{cases}
\end{equation}
这是把公式
\begin{align}
  \omega(e^{k\theta J}v,e^{\ell\theta J}Jw)
  ={}&\cos(k\theta)\cos(\ell\theta)\omega(v,Jw)
  +\sin(k\theta)\sin(\ell\theta)\omega(w,Jv)\notag\\
  &+\sin(k\theta)\cos(\ell\theta)\omega(Jv,Jw)
  +\cos(k\theta)\sin(\ell\theta)\omega(w,v)\notag\\
  ={}&\cos((k-\ell)\theta)\langle v,w\rangle
  +\sin((k-\ell)\theta)\omega(v,w)
  \label{eq:4-4-8}
\end{align}
在区间 $0\le\theta\le2\pi$ 上积分得到的。（最后一个等号用到了 $J$ 与 $\omega$
\BookRecordIndex{T0348}相容。若 $J$ 只被 $\omega$ \BookRecordIndex{T0372}驯服，则当 $k+\ell\ne0$ 时，式~\eqref{eq:4-4-7}
仍然成立，但当 $k+\ell=0$ 时并不成立。）现在把 $\gamma$ 写成 Fourier 级数
\begin{equation}\label{eq:4-4-9}
  \gamma(\theta)=\sum_{k=-\infty}^{\infty}e^{k\theta J}v_k.
\end{equation}
于是由式~\eqref{eq:4-4-7}，
\[
  \begin{aligned}
  A(\gamma)
  &=-\frac12\int_0^{2\pi}
    \omega(\gamma(\theta),\dot\gamma(\theta))\,d\theta\\
  &=-\frac12\sum_{k,\ell}\ell\int_0^{2\pi}
    \omega(e^{k\theta J}v_k,e^{\ell\theta J}Jv_\ell)\,d\theta\\
  &=-\pi\sum_k k|v_k|^2,
  \end{aligned}
\]
并且
\[
  \begin{aligned}
  E(\gamma)
  &=\frac12\int_0^{2\pi}
    \omega(\dot\gamma(\theta),J\dot\gamma(\theta))\,d\theta\\
  &=\frac12\sum_{k,\ell}k\ell\int_0^{2\pi}
    \omega(e^{k\theta J}v_k,e^{\ell\theta J}Jv_\ell)\,d\theta\\
  &=\pi\sum_k k^2|v_k|^2.
  \end{aligned}
\]
这就\BookRecordIndex{T0326}推出式~\eqref{eq:4-4-6}。

为由式~\eqref{eq:4-4-6} \BookRecordIndex{T0326}推出式~\eqref{eq:4-4-4}，先设 $\gamma$ 是浸入，并以弧长
重新参数化。也就是说，假设每个 $\theta$ 都满足 $\dot\gamma(\theta)\ne0$，并选取
重参数化 $\alpha:\mathbb R\to\mathbb R$，使其提升满足正确的拟周期条件
\sourcefootnote{原书此处写作 $\alpha(\theta+2\pi)=\alpha(\theta)$；这会使
$\alpha$ 不可能成为圆周的提升微分同胚。由闭路参数的周期性，唯一应改为
$\alpha(\theta+2\pi)=\alpha(\theta)+2\pi$。}：
\[
  \alpha(\theta+2\pi)=\alpha(\theta)+2\pi
\]
以及
\[
  |\partial_\theta(\gamma\circ\alpha)|\equiv\frac{L(\gamma)}{2\pi}.
\]
于是由式~\eqref{eq:4-4-6}，
\[
  |A(\gamma)|=|A(\gamma\circ\alpha)|
  \le E(\gamma\circ\alpha)
  =\frac1{4\pi}L(\gamma)^2.
\]
这对浸入曲线证明了式~\eqref{eq:4-4-4}。为证明一般情形，选取趋于零的向量列
$a_\nu\in V$，使得对每个 $\theta$ 和每个 $\nu$，
\[
  a_\nu\ne e^{-\theta J}J\dot\gamma(\theta).
\]
那么
\[
  \gamma_\nu(\theta):=\gamma(\theta)+e^{\theta J}a_\nu
\]
是一列以 $C^\infty$ 拓扑收敛到 $\gamma$ 的浸入。因而对每个 $\nu$，
$|A(\gamma_\nu)|\le L(\gamma_\nu)^2/(4\pi)$；令 $\nu\to\infty$，即得 $\gamma$
的结论。这证明了式~\eqref{eq:4-4-4}。

为证明式~\eqref{eq:4-4-5}，注意由式~\eqref{eq:4-4-8}，
\begin{equation}\label{eq:4-4-10}
  \int_0^\pi\omega(e^{k\theta J}v,e^{\ell\theta J}Jw)\,d\theta
  =
  \begin{cases}
    \pi\langle v,w\rangle,&k=\ell,\\
    0,&k\ne\ell,
  \end{cases}
  \qquad v,w\in\Lambda.
\end{equation}
现在设 $\gamma:[0,\pi]\to V$ 是满足
$\gamma(0),\gamma(\pi)\in\Lambda$ 的光滑道路，并把 $\gamma$ 写成
式~\eqref{eq:4-4-9} 的 Fourier 级数，其中 $v_k\in\Lambda$。由
式~\eqref{eq:4-4-10}，
\[
  A(\gamma)=-\frac\pi2\sum_k k|v_k|^2,
  \qquad
  E(\gamma)=\frac\pi2\sum_k k^2|v_k|^2.
\]
所以在时间区间 $[0,\pi]$ 上的拉格朗日情形，式~\eqref{eq:4-4-6} 仍然成立。对浸入
曲线，重参数化便由式~\eqref{eq:4-4-6} 得到估计~\eqref{eq:4-4-5}。一般情形则在上述
逼近论证中取 $a_\nu\in\Lambda$。这就证明了引理~\ref{lem:4-4-4}。
\end{proof}

下面的例子来自 Pavel Giterman。它说明，当 $\dim V\ge4$ 时，被 $\omega$ \BookRecordIndex{T0372}驯服的
复结构之\BookRecordIndex{T0099}等周不等式不可能以独立于 $J$ 的常数成立。（本书 2004 年版曾声称：对被
$\omega$ \BookRecordIndex{T0372}驯服的复结构 $J$，估计~\eqref{eq:4-4-4} 与~\eqref{eq:4-4-5} 都成立；
Joseph Bernstein 指出了这个错误。）

\begin{example}\label{ex:4-4-5}
把 $V=\mathbb R^{2n}$ 的元素写成 $z=(x,y)$，其中 $x,y\in\mathbb R^n$。以
\[
  \omega:=\sum_i dx_i\wedge dy_i
\]
表示标准\BookRecordIndex{T0363}辛形式，并任取 $Q\in\operatorname{GL}(n,\mathbb R)$。复结构
\[
  J_Q:=
  \begin{pmatrix}
    0&-Q^{-1}\\
    Q&0
  \end{pmatrix}
\]
是否被 $\omega$ \BookRecordIndex{T0372}驯服取决于矩阵的正定性条件\sourcefootnote{原书断言它对任意
$Q\in\operatorname{GL}(n,\mathbb R)$ 都被 $\omega$ \BookRecordIndex{T0372}驯服，但一般并不成立，例如
$Q=-I$ 即给出反例；相应二次型为正定才有\BookRecordIndex{T0372}驯服性，
等价地须要求 $Q$ 的对称部分正定。下述具体矩阵族在 $0<\varepsilon\le1$ 时对称部分为
$\varepsilon I$，故后续例族不受这一一般性过强的断言影响。}。$\mathbb R^{2n}$ 上相应的
二次型为
\[
  |z|_Q^2:=\omega(z,J_Qz)
  =\langle x,Qx\rangle+\langle y,Q^{-1}y\rangle.
\]
取 $n=2$、$0<\varepsilon\le1$，并选择
\[
  Q:=
  \begin{pmatrix}
    \varepsilon&-\sqrt{1-\varepsilon^2}\\
    \sqrt{1-\varepsilon^2}&\varepsilon
  \end{pmatrix}
  =(Q^{-1})^T.
\]
于是 $|z|_Q^2=\varepsilon|z|^2$，故把式~\eqref{eq:4-4-4} 中 $V$ 上的范数换成
$|\cdot|_Q$ 后，$c_Q:=1/(4\pi\varepsilon)$ 是\BookRecordIndex{T0099}等周不等式的最佳常数。
\end{example}

\begin{proof}[定理~\ref{thm:4-4-1} 的证明]
先假设 $J$ 与 $\omega$ \BookRecordIndex{T0348}相容。用有限多个 \BookRecordIndex{T0081}达布坐标卡
$\phi:U\to\mathbb R^{2n}$ 覆盖 $M$，并使每个坐标卡都能光滑延拓到 $\overline U$
的一个开邻域。现在选取 $\delta>0$，使得对每个 $x_0\in M$，上述有限族中都有一个
\BookRecordIndex{T0081}达布坐标卡 $\phi$ 满足：
\begin{enumerate}[label=(\roman*),leftmargin=3em]
\item 球 $B_\delta(x_0)$ 包含在 $U$ 中；
\item 对每个 $x\in B_\delta(x_0)$，
\[
  \lVert d\phi(x)\rVert^2
  \bigl\lVert\phi_*J(\phi(x))-\phi_*J(\phi(x_0))\bigr\rVert
  \le\varepsilon:=4\pi c-1.
\]
\end{enumerate}
这里，$\lVert d\phi(x)\rVert$ 表示线性映射
$d\phi(x):T_xM\to\mathbb R^{2n}$ 关于 $T_xM$ 上式~\eqref{eq:2-1-1} 的度量和
$\mathbb R^{2n}$ 上欧氏范数的算子范数；第二项
$\lVert\phi_*J(\phi(x))-\phi_*J(\phi(x_0))\rVert$ 则表示
$\mathbb R^{2n\times2n}$ 上的欧氏矩阵范数。

设 $\gamma:S^1\to M$ 是长度 $\ell(\gamma)<\delta$ 的闭路，并选取上述 Darboux
坐标卡 $\phi:U\to\mathbb R^{2n}$，使它包含点 $x_0:=\gamma(0)$。于是
$\gamma(S^1)\subset B_\delta(x_0)\subset U$。令
\[
  J_0:=\phi_*J(\phi(x_0))
  =d\phi(x_0)J(x_0)d\phi(x_0)^{-1},
\]
这是 $J(x_0)$ 在 $\mathbb R^{2n}$ 上诱导的与 $\omega_0$ \BookRecordIndex{T0348}相容的复结构。对每个
$x\in B_\delta(x_0)$ 和每个 $\xi\in T_xM$，都有
\[
  \begin{aligned}
  |d\phi(x)\xi|_{J_0}^2
  ={}&|\xi|_J^2
  +\omega_0\bigl(d\phi(x)\xi,
    (J_0-d\phi(x)J(x)d\phi(x)^{-1})d\phi(x)\xi\bigr)\\
  \le{}&\left(1+\lVert d\phi(x)\rVert^2
    \bigl\lVert\phi_*J(\phi(x_0))-\phi_*J(\phi(x))\bigr\rVert\right)|\xi|_J^2\\
  \le{}&(1+\varepsilon)|\xi|_J^2.
  \end{aligned}
\]
所以由引理~\ref{lem:4-4-4}，
\[
  |a(\gamma)|=|A(\phi\circ\gamma)|
  \le\frac1{4\pi}L(\phi\circ\gamma)^2
  \le\frac{1+\varepsilon}{4\pi}\ell(\gamma)^2
  =c\ell(\gamma)^2.
\]
这里 $L(\phi\circ\gamma)$ 关于 $\mathbb R^{2n}$ 上由 $\omega_0$ 和 $J_0$ 诱导的度量
理解。这证明了 (i)。

(ii) 的证明完全相同，只是现在必须选取满足
$\phi(U\cap L)=\phi(U)\cap\mathbb R^n$ 的 \BookRecordIndex{T0081}达布坐标卡，并令
$\varepsilon:=2\pi c-1$。为证明 (iii)，选取一个与 $\omega$ \BookRecordIndex{T0348}相容的\BookRecordIndex{T0088}殆复结构 $J_0$。
于是 (i) 与 (ii) 对 $J_0$ 成立。由于 $M$ 紧，道路 $\gamma$ 关于 $J_0$ 与 $J$ 的
长度满足
\[
  \frac{\ell_0(\gamma)}{c_0}
  \le\ell(\gamma)\le c_0\ell_0(\gamma),
\]
其中常数 $c_0$ 与 $\gamma$ 无关。这就证明了定理~\ref{thm:4-4-1}。
\end{proof}

\begin{exercise}\label{ex:4-4-6}
证明：若 $(M,J)$ 是配备黎曼度量 $g$ 的紧殆复流形，则存在正常数 $\delta,c$，使每条
长度小于 $\delta$ 的闭曲线 $\gamma$ 都是某个圆盘 $u:\mathbb D\to M$ 的边界，并且
\[
  E(u)<c\ell(\gamma)^2.
\]
\end{exercise}
